\documentclass[11pt,a4paper]{article}

\usepackage[T1]{fontenc}
\usepackage[utf8]{inputenc}
\usepackage[english]{babel}
\usepackage{lmodern}
\usepackage{microtype}
\usepackage[a4paper,margin=20mm,headheight=15pt]{geometry}
\usepackage{graphicx}
\usepackage{booktabs}
\usepackage{tabularx}
\usepackage{array}
\usepackage{xcolor}
\usepackage{amsmath}
\usepackage{caption}
\usepackage{float}
\usepackage{enumitem}
\usepackage{fancyhdr}
\usepackage[hidelinks]{hyperref}
\usepackage{xurl}

\graphicspath{{figures/}}
\definecolor{nativeblue}{HTML}{2563EB}
\definecolor{flexorange}{HTML}{EA580C}
\definecolor{passivepurple}{HTML}{7C3AED}
\definecolor{rtkgreen}{HTML}{16A34A}
\definecolor{slate}{HTML}{334155}
\definecolor{lightgray}{HTML}{F1F5F9}
\definecolor{warning}{HTML}{B45309}

\newcommand{\native}{\textcolor{nativeblue}{Native DS-TWR}}
\newcommand{\passive}{\textcolor{passivepurple}{Passive DS-TWR}}
\newcommand{\flex}{\textcolor{flexorange}{FlexTDOA}}
\newcommand{\rtk}{\textcolor{rtkgreen}{GPS RTK}}
\newcommand{\code}[1]{\texttt{\detokenize{#1}}}
\newcolumntype{Y}{>{\raggedright\arraybackslash}X}

\pagestyle{fancy}
\fancyhf{}
\lhead{\small Four-method field comparison}
\rhead{\small August 3, 2026}
\cfoot{\thepage}

\hypersetup{
  pdftitle={Field comparison of FlexTDOA, Native DS-TWR, Passive DS-TWR and GPS RTK},
  pdfauthor={UWB ESP-IDF technical report},
  pdfsubject={Static outdoor field comparison using a common RTK-derived reference frame},
  pdfkeywords={UWB, DW3000, FlexTDOA, DS-TWR, Passive DS-TWR, GNSS, GPS, RTK}
}

\begin{document}
\pagenumbering{gobble}

\begin{titlepage}
  \centering
  \vspace*{14mm}
  {\Huge\bfseries Field comparison of four positioning methods\par}
  \vspace{4mm}
  {\Large \flex{}, \native{}, \passive{} and \rtk{}\par}
  \vspace{7mm}
  {\large Static outdoor data, common RTK-derived ENU frame, no UWB filtering\par}
  \vspace{15mm}

  \begin{tabular}{rl}
    Test date: & August 3, 2026, 16:58--17:14 EEST\\
    Platform: & 5 ESP32-S3 + DW3000 + PX1105R-01 modules, Raspberry Pi 5\\
    Firmware: & \texttt{8fce48f-dirty}; repository HEAD \texttt{0c922dd}\\
    Configuration: & T1 tag; A2, A3, A4, A5 anchors\\
    Acquisition: & three consecutive 300 s protocol blocks; GNSS configured at 1 Hz\\
    Environment: & outdoor, static geometry, no temporal filtering\\
  \end{tabular}

  \vfill
  \colorbox{lightgray}{%
    \begin{minipage}{0.89\textwidth}
      \vspace{3mm}
      \textbf{Preliminary result.}
      Under this single static field geometry, \native{} produced the tightest
      unfiltered position cloud: \textbf{1.39 cm CEP50} and
      \textbf{2.99 cm P95}. \passive{} delivered
      \textbf{26.12 complete positions/s}, approximately 11.4 times the Native
      rate, with \textbf{6.57 cm CEP50} and \textbf{10.06 cm P95}.
      The embedded \flex{} path published \textbf{54.68 positions/s}, while
      the common complete-slot reconstruction yielded 16.97 positions/s.
      Relative to the selected RTK reference, the three UWB RMSE values were
      13.88 cm, 11.52 cm, and 10.15 cm, respectively.
      \vspace{3mm}
    \end{minipage}}

  \vfill
  \colorbox{warning!10}{%
    \begin{minipage}{0.89\textwidth}
      \vspace{2mm}
      \textbf{RTK integrity warning.}
      A3 reported RTK fixed throughout most of the campaign but changed by
      \textbf{2.62 m} between protocol blocks. The report therefore uses the
      only RTK block whose complete six-pair anchor geometry agrees with the
      static UWB anchor-range consensus. RTK is a valuable common reference
      here, but not an independently surveyed ground truth.
      \vspace{2mm}
    \end{minipage}}
  \vspace{10mm}
  {\large Reproducible preliminary technical report\par}
\end{titlepage}

\pagenumbering{roman}
\tableofcontents
\clearpage
\pagenumbering{arabic}

\section{Executive summary}

This campaign is the first common data set containing all three implemented
UWB protocols and GPS RTK. Each UWB protocol ran for five minutes in the same
static outdoor layout. GNSS remained active on all modules, nominally at
1 Hz. The comparison deliberately separates two questions:

\begin{enumerate}[leftmargin=7mm]
  \item \textbf{Accuracy against the RTK reference}: horizontal distance
  between a UWB solution and the robust RTK-fixed T1 position.
  \item \textbf{Precision}: radial dispersion about each method's own robust
  median, which does not reward or penalize a constant offset.
\end{enumerate}

All three UWB raw streams were reconstructed with the same unweighted,
unfiltered two-dimensional offline solver and the same anchor coordinates.
This avoids comparing different dashboard filters or embedded solver policies.
The ESP32-published streams are listed separately as end-to-end results.

\begin{table}[H]
  \centering
  \caption{Primary common-solver results. Accuracy is relative to the selected
  RTK reference; precision is dispersion about each method's own median.}
  \footnotesize
  \setlength{\tabcolsep}{4pt}
  \begin{tabular}{lrrrrrr}
    \toprule
    Method & $N$ & Hz & Accuracy RMSE & Accuracy P95 & Precision CEP50 & Precision P95\\
    \midrule
    \flex{} & 5,093 & 16.97 & 13.88 cm & 23.15 cm & 9.23 cm & 19.64 cm\\
    \native{} & 688 & 2.29 & 11.52 cm & 13.50 cm & 1.39 cm & 2.99 cm\\
    \passive{} & 7,838 & 26.12 & 10.15 cm & 14.97 cm & 6.57 cm & 10.06 cm\\
    \rtk{} & 544 & 0.60 observed & reference & reference & 0.50 cm & 1.10 cm\\
    \bottomrule
  \end{tabular}
\end{table}

The result is nuanced. Native DS-TWR is clearly the most repeatable UWB method
in this static trial. Passive DS-TWR happens to have the lowest RMSE against
the chosen RTK position, but its wider cloud shows that this is partly a
favorable cancellation of systematic bias; it is not more precise than Native
DS-TWR. FlexTDOA remains the fastest end-to-end implementation, while Passive
DS-TWR is the fastest source of complete frames reconstructed by the common
solver in this particular capture.

\section{System and tested protocol settings}

\subsection{Hardware and radio configuration}

The system contains one receive-capable tag (T1) and four anchors (A2--A5).
Every node combines an ESP32-S3, a Qorvo DW3000 UWB transceiver, and a SkyTraq
PX1105R-01 GNSS receiver. All UWB methods used channel 9. GNSS RTK corrections
were distributed through Wi-Fi. The tag and anchors were not moved between the
three acquisition blocks.

\begin{table}[H]
  \centering
  \caption{Runtime parameters captured in the metadata.}
  \small
  \begin{tabularx}{\textwidth}{p{34mm}p{31mm}Y}
    \toprule
    Method & Frame/slot & Relevant settings\\
    \midrule
    \flex{} & 3 ms slot, 12 ms four-anchor frame
      & 250 $\mu$s request subslot; 1,000 $\mu$s request processing;
        250 $\mu$s response subslot and processing; 250 $\mu$s guard time\\
    \native{} & 100 ms per anchor; 410 ms nominal four-anchor frame
      & 10 ms frame gap; 90 ms RX timeout; 20 ms response delay;
        20 ms final delay; native POLL--RESP--FINAL exchange\\
    \passive{} & 5 ms exchange slot; 1 ms gap
      & Robust Rotating schedule; 4 ms RX timeout; 1,000 $\mu$s response;
        1,000 $\mu$s final delay; receive-only tag\\
    \rtk{} & configured 1 Hz
      & PX1105R-01, NTRIP corrections through Wi-Fi; fixed solutions retained\\
    \bottomrule
  \end{tabularx}
\end{table}

\subsection{Position observables}

Native DS-TWR measures four absolute tag-to-anchor distances. The position is
obtained by unweighted 2-D trilateration. FlexTDOA and Passive DS-TWR produce
range-difference observations of the form

\begin{equation}
 z_{ij} \simeq \lVert\mathbf p-\mathbf a_j\rVert
             -\lVert\mathbf p-\mathbf a_i\rVert,
\end{equation}

and the common solver minimizes

\begin{equation}
 \min_{\mathbf p}\sum_{(i,j)}
 \left(\lVert\mathbf p-\mathbf a_j\rVert
 -\lVert\mathbf p-\mathbf a_i\rVert-z_{ij}\right)^2.
\end{equation}

No median filter, moving average, Kalman filter, or trajectory smoothing was
applied to any UWB position in the primary comparison. Physical corrections
inside the radio measurements are not treated as temporal filtering.

\section{Acquisition and common analysis method}

\subsection{Data blocks}

The sequence was FlexTDOA, Native DS-TWR, then Passive DS-TWR. Each acquisition
lasted 300 s after a 15 s warm-up. The collector polled module status at 50 Hz
and requested GNSS at 1 Hz. Data were deduplicated by their native protocol
keys before analysis.

\begin{table}[H]
  \centering
  \caption{Captured data volume after key-based deduplication.}
  \small
  \begin{tabular}{lrrrr}
    \toprule
    Block & UWB observations & Anchor ranges & ESP32 positions & GNSS events\\
    \midrule
    FlexTDOA & 192,787 & 203,581 & 16,408 & 924\\
    Native DS-TWR & 2,887 ranges & 176 & 0 & 937\\
    Passive DS-TWR & 37,320 & 44,728 & 7,839 & 932\\
    \bottomrule
  \end{tabular}
\end{table}

The absence of Native embedded positions is intentional: during this capture
the Native position solver was running on the Raspberry Pi. Its raw ranges are
therefore reconstructed offline by the common solver.

\subsection{Common ENU frame and accuracy reference}

RTK-fixed WGS84 latitude, longitude, and altitude were converted to ECEF and
then to a local East--North--Up frame whose origin is the robust A2 position.
The robust T1 median across all blocks is the horizontal accuracy reference.
The tag changed by only 1.5 cm between block centers, so pooling its fixed
samples is appropriate for this static test.

The anchor geometry could not be pooled blindly: A3 entered mutually
inconsistent RTK-fixed solutions. Section~\ref{sec:rtk-integrity} documents the
selection of the coherent geometry. The resulting geometry is shown below.

\begin{figure}[H]
  \centering
  \includegraphics[width=0.76\textwidth]{01_rtk_geometry.pdf}
  \caption{Common analysis geometry derived from the coherent RTK block. The
  T1 marker is the robust RTK-fixed reference position.}
\end{figure}

\subsection{Frame assembly}

Native DS-TWR uses a complete set only when all four ranges carry the same
reconstructed cycle. FlexTDOA uses three responder observations belonging to
one native slot. Passive DS-TWR emits the three responder exchanges in three
consecutive slots; the common frame identifier is recovered from the ESP32
slot identifier and responder index. No Raspberry Pi wall-clock synchronization
is used to number UWB frames.

\section{RTK integrity audit}\label{sec:rtk-integrity}

The GNSS receiver reported RTK fixed for more than 96\% of the retained samples
on every node. Nevertheless, solution status alone was insufficient: the A3
block center moved by 2.62 m while A2, A4, A5, and T1 each remained within
1.8 cm across the campaign.

\begin{table}[H]
  \centering
  \caption{Maximum three-dimensional displacement between consecutive RTK
  block centers.}
  \begin{tabular}{lrrrrr}
    \toprule
    Module & T1 & A2 & A3 & A4 & A5\\
    \midrule
    Maximum span & 1.47 cm & 1.21 cm & \textbf{2.62 m} & 1.70 cm & 1.48 cm\\
    \bottomrule
  \end{tabular}
\end{table}

Because the physical anchor layout was static, the six robust anchor-to-anchor
UWB medians provide a cross-sensor integrity check. For each RTK epoch, the
six RTK 3-D separations were compared with the corresponding UWB consensus.

\begin{table}[H]
  \centering
  \caption{RTK anchor-geometry integrity against the six-pair static UWB
  consensus. The selected geometry is bold.}
  \begin{tabular}{lrrr}
    \toprule
    RTK block & Complete pairs & Pair RMSE & Maximum residual\\
    \midrule
    FlexTDOA block & 6 & 1.068 m & 2.232 m\\
    Native DS-TWR block & 6 & 0.848 m & 1.757 m\\
    \textbf{Passive DS-TWR block} & \textbf{6} & \textbf{0.071 m} & \textbf{0.137 m}\\
    \bottomrule
  \end{tabular}
\end{table}

The Passive block is the only complete RTK geometry compatible with the radio
distances and is used for all three UWB solvers. This procedure prevents a
known false-fixed coordinate from generating metre-scale fictitious UWB
errors. It also creates a limited dependence between the geometry reference
and UWB, so this report must not be interpreted as a metrological calibration
against independent survey equipment.

\begin{figure}[H]
  \centering
  \includegraphics[width=0.94\textwidth]{08_rtk_quality.pdf}
  \caption{RTK status availability and T1 reference-stream statistics. The
  high fixed percentage does not reveal the A3 integrity failure.}
\end{figure}

\section{Position results}

\subsection{Accuracy relative to RTK}

\begin{figure}[H]
  \centering
  \includegraphics[width=0.96\textwidth]{02_position_clouds.pdf}
  \caption{Unfiltered position clouds in the same RTK-derived coordinate
  frame. The black cross is the RTK-fixed T1 median. Identical axes make the
  different dispersions directly visible.}
\end{figure}

\begin{figure}[H]
  \centering
  \includegraphics[width=0.88\textwidth]{03_error_cdf.pdf}
  \caption{CDF of horizontal UWB error relative to T1 RTK. The green curve is
  RTK self-dispersion, not an independent RTK accuracy measurement.}
\end{figure}

The Passive DS-TWR median lies closer to T1 RTK than the Native DS-TWR median,
which explains its slightly lower RMSE. Native DS-TWR remains much tighter:
its complete error distribution spans only 10.0 cm from the observed minimum
to maximum, whereas Passive contains a 51.4 cm maximum error. Therefore the
accuracy ranking at this one point must be read together with precision and
range bias, not from RMSE alone.

\subsection{Static precision}

\begin{figure}[H]
  \centering
  \includegraphics[width=0.88\textwidth]{04_precision_boxplot.pdf}
  \caption{Radial displacement from each method's own median. Plot outliers
  are hidden for readability but retained in all tabulated metrics.}
\end{figure}

Native DS-TWR is the clear static-precision winner: 1.39 cm CEP50 and 2.99 cm
P95. GPS RTK is approximately three times tighter at P95, while Passive DS-TWR
and FlexTDOA are respectively 3.4 and 6.6 times wider than Native at P95.

\begin{figure}[H]
  \centering
  \includegraphics[width=0.95\textwidth]{05_error_time.pdf}
  \caption{Horizontal UWB error versus time. No temporal filtering or
  ground-truth-based rejection was applied.}
\end{figure}

\section{Rate, delivery, and end-to-end output}

\begin{figure}[H]
  \centering
  \includegraphics[width=0.94\textwidth]{06_rate_delivery.pdf}
  \caption{Common-solver update rate and completeness. Candidate-unit
  definitions are protocol-specific, so completeness should be interpreted
  as a delivery diagnostic rather than a direct efficiency ranking.}
\end{figure}

\begin{table}[H]
  \centering
  \caption{Frame assembly and delivery diagnostics.}
  \small
  \begin{tabular}{lrrrrr}
    \toprule
    Method & Candidates & Complete & Complete & Median interval & P95 interval\\
    \midrule
    \flex{} & 97,543 & 5,093 & 5.22\% & 48.3 ms & 193.9 ms\\
    \native{} & 735 & 688 & 93.61\% & 410.3 ms & 811.2 ms\\
    \passive{} & 16,013 & 7,838 & 48.95\% & 29.9 ms & 134.0 ms\\
    \bottomrule
  \end{tabular}
\end{table}

FlexTDOA's primary common-solver row deliberately requires all three responses
from one physical slot. The product's embedded solver publishes more often by
using the current valid observation state. This produces the following
end-to-end precision result, which is useful operationally but is not identical
to the complete-frame comparison above.

\begin{table}[H]
  \centering
  \caption{Positions published by the running ESP32 solver. Native used the
  Raspberry solver during this capture and therefore has no embedded row.}
  \begin{tabular}{lrrrr}
    \toprule
    Method & $N$ & Published Hz & Precision CEP50 & Precision P95\\
    \midrule
    \flex{} embedded & 16,408 & 54.68 & 3.44 cm & 7.51 cm\\
    \passive{} embedded & 7,839 & 26.13 & 5.06 cm & 8.94 cm\\
    \bottomrule
  \end{tabular}
\end{table}

This distinction resolves an apparent contradiction: FlexTDOA is the fastest
live implementation, but Passive DS-TWR produced more fully coherent raw
frames for the strict common solver in this capture.

\section{Anchor-range consistency}

\begin{figure}[H]
  \centering
  \includegraphics[width=0.96\textwidth]{07_anchor_range_bias.pdf}
  \caption{Median UWB anchor range minus the coherent RTK 3-D separation.
  These biases include both UWB hardware bias and residual RTK-geometry error.}
\end{figure}

All three protocols show the same sign pattern on most pairs, suggesting that
a substantial component is hardware/path dependent rather than solver noise.
The largest common discrepancy is A4--A5, approximately $-13.4$ to $-16.4$ cm.
Native DS-TWR has the smallest range dispersion on most pairs, generally
1.3--3.1 cm standard deviation; Passive DS-TWR is similar on five pairs, while
FlexTDOA anchor-range telemetry is wider at about 4.6--5.8 cm.

This pair structure is a clear calibration opportunity. A future independent
survey should estimate antenna-delay and pair biases without using RTK epochs
that fail the geometry integrity check.

\section{Conclusions}

\begin{enumerate}[leftmargin=7mm]
  \item \textbf{Native DS-TWR is the strongest static precision reference.}
  It achieves 2.99 cm P95 with no temporal filtering, but only 2.29 complete
  positions/s under the conservative 410 ms four-anchor profile.

  \item \textbf{Passive DS-TWR is the best complete-frame speed/precision
  compromise in this capture.} It reaches 26.12 positions/s, 6.57 cm CEP50,
  and 10.06 cm P95 while the tag remains receive-only.

  \item \textbf{FlexTDOA remains the fastest live output.} The ESP32 publishes
  54.68 positions/s with 7.51 cm P95 self-dispersion. Its strict, complete-slot
  common reconstruction is slower and noisier because only 5.22\% of candidate
  slots contain every required observation in the captured stream.

  \item \textbf{Accuracy and precision lead to different rankings.} Passive
  DS-TWR has the lowest RMSE against this RTK reference, but Native DS-TWR has
  by far the smallest cloud. A single favorable bias cancellation must not be
  mistaken for general accuracy superiority.

  \item \textbf{RTK fixed is not an integrity guarantee.} A3 produced a
  2.62 m false-fixed displacement. Cross-checking static inter-anchor geometry
  is mandatory before RTK is used as an accuracy reference.
\end{enumerate}

\paragraph{Recommended next experiment.}
Repeat randomized five-minute blocks at several tag positions and include a
dynamic trajectory. Survey anchor coordinates independently or use a verified
local reference, and randomize protocol order to separate protocol behavior
from time-varying RF and GNSS conditions. Retain both common complete-frame and
actual end-to-end results.

\section{Limitations}

This report describes one static geometry and three sequential five-minute
blocks. It does not test edge geometry, NLOS, antenna orientation, speed, or
multi-tag scaling. The protocol order was not randomized. GNSS was configured
at 1 Hz, but the distinct retained T1 fixed stream was observed at 0.60 Hz over
the full 900 s campaign. Most importantly, RTK is part of the tested system and
the coherent anchor epoch was selected using static UWB pair consistency; the
reported UWB ``accuracy'' is therefore a practical cross-method comparison,
not traceable absolute accuracy.

\section{Reproducibility}

The report directory contains:

\begin{itemize}[leftmargin=7mm]
  \item raw JSONL captures and per-block metadata under \code{data/};
  \item \path{analysis_summary.json}, \path{position_metrics.csv},
    \path{anchor_pair_metrics.csv}, and \path{rtk_geometry_integrity.csv};
  \item all common-solver positions in \path{common_solver_positions.csv};
  \item vector and raster figures under \code{figures/};
  \item checksums in \code{SHA256SUMS}.
\end{itemize}

The deterministic analysis command is:

\begin{quote}
\small\ttfamily
python3 tools/uwb\_gps\_rtk\_report.py \textbackslash\\
\hspace*{2em}reports/uwb\_gps\_rtk\_comparison\_20260803
\end{quote}

\end{document}
